\section*{الفصل \ref{ch:b2:blood-pressure} --- تنظيم ضغط الدم والجهد}

\begin{solution}{exo:b2:blood-pressure:1}
$\bar P - P_{\text{وريدي}} = Q\,R = f\,V_s\,R$. معدل \omterm{def:b2:heart:anatomy}{القلب} (أي
\omterm{prop:b2:heart:conduction}{العقدة الجيبية الأذينية}، تحت المبهم والأعصاب الودّية)؛ \omterm{prop:b2:heart:cycle}{وحجم
الضربة} (أي \omterm{def:b2:heart:anatomy}{البطين}: امتلاؤه الذي تقرره الأوردة، وقابليته
للانقباض التي تقررها الأعصاب الودّية)؛ والمقاومة (أي
العضل الأملس الشريني، تحت الأعصاب الودّية والهرمونات
ونواتج الاستقلاب الموضعية).
\end{solution}

\begin{solution}{exo:b2:blood-pressure:2}
المستشعرات: مستقبلات التمدد في الجيبين السباتيين وقوس الأبهر.
والمركز: المراكز القلبية الوعائية في البصلة السيسائية. والمنفِّذات:
المبهم (المعدل) والأعصاب الودّية (المعدل والقابلية للانقباض
والمقوية الشرينية والوريدية). وتغذية راجعة سالبة، بتأخر نحو ثانية.
ودون المستشعرات يبقى الضغط المتوسط على مدى يوم
سويًا لكن الضغط من دقيقة إلى دقيقة يتأرجح بعنف مع كل
وضعة وانفعال.
\end{solution}

\begin{solution}{exo:b2:blood-pressure:3}
تطلق الكلية (خلايا شريناتها) الرينين حين تهبط ترويتها أو
الصوديوم المسلَّم إليها أو تنطلق أعصابها الودّية؛ ويقطع الرينين
مولد الأنجيوتنسين (من الكبد) إلى أنجيوتنسين I؛ ويحوّله إنزيم الرئة إلى
أنجيوتنسين II، الذي يقبض \omterm{def:b2:blood-circulation:vessels}{الشرينات} ويحرض
قشر الكظر على إفراز الألدوستيرون؛ والألدوستيرون يجعل الكلية
تحتبس الصوديوم والماء، فيرتفع حجم الدم.
\end{solution}

\begin{solution}{exo:b2:blood-pressure:4}
التوسع الاستقلابي \omterm{def:b2:blood-circulation:vessels}{لشرينات} العضلة نفسها بنواتج
استقلابها؛ \omterm{thm:b2:heart:starling}{ونتاج قلبي} بخمسة أضعاف تسوقه الأعصاب
الودّية والمضخة العضلية؛ والتقبض الودّي \omterm{def:b2:blood-circulation:vessels}{لشرينات}
الأمعاء والكليتين والعضلات المستريحة، وهي التي تتنازل عن
جريانها.
\end{solution}

\begin{solution}{exo:b2:blood-pressure:5}
$83\times 1.2 = \qty{100}{mmHg}$؛ و$333\times 0.35 = \qty{117}{mmHg}$.
وهبوط المقاومة $1.2/0.35 = 3.4$: فيرتفع نصف القطر $3.4^{1/4} =
1.36$.
\end{solution}

\begin{solution}{exo:b2:blood-pressure:6}
$30/(1 + G)$: أي 10 و6 و\qty{3.3}{mmHg}. وبنصف الربح يتضاعف
الهبوط المتبقي: فتدوخ عند الوقوف، وترى رماديًا، وقد
يُغشى عليها.
\end{solution}

\begin{solution}{exo:b2:blood-pressure:7}
$\rho g h$: أي $1060\times 9.81\times 0.4 = \qty{4160}{Pa} =
\qty{31}{mmHg}$؛ و$1060\times 9.81\times 1.2 = \qty{12500}{Pa} =
\qty{94}{mmHg}$. فالدماغ $95 - 31 = \qty{64}{mmHg}$؛ والكاحل $95 + 94 =
\qty{189}{mmHg}$. ورأس الزرافة على \qty{2.5}{m} يكلف
\qty{195}{mmHg}: فهي تحتاج ضغطًا متوسطًا نحو \qty{250}{mmHg} عند
\omterm{def:b2:heart:anatomy}{القلب}، وجدر أوعية ثخينة وجلدًا مشدودًا في الساقين.
\end{solution}

\begin{solution}{exo:b2:blood-pressure:8}
$Q = 280/(190 - 130) = \qty{4.7}{L/min}$. وعند الرياضية: $5000/(200 - 30) =
\qty{29}{L/min}$؛ \omterm{prop:b2:heart:cycle}{وحجم الضربة} $\num{29400}/190 = \qty{155}{mL}$.
\end{solution}

\begin{solution}{exo:b2:blood-pressure:9}
$35/5 = \qty{7}{mmHg}$ هبوطًا: أي نحو \qty{88}{mmHg}. وإعادة
الحجم: امتصاص سائل الخلال إلى الشعيرات
(من دقائق إلى ساعة)؛ وتقليل البول تحت \omterm{prop:b2:blood-pressure:raas}{الهرمون المضاد للإبالة}
والألدوستيرون (من ساعات إلى أيام)؛ والعطش والشرب (ساعات)؛
والكريات الحمر من النقي (أسابيع). والجلد شاحب بارد لأن
المنعكس أغلق شريناته ليحفظ الضغط للدماغ
\omterm{def:b2:heart:anatomy}{والقلب}.
\end{solution}

\begin{solution}{exo:b2:blood-pressure:10}
العضلة عند $R = 0.5$: $95/0.5 = \qty{190}{mL/s}$، أي \qty{11.4}{L/min}؛
ومع \qty{4}{L/min} لبقية الأعضاء يحتاج \omterm{def:b2:heart:anatomy}{القلب}
\qty{15.4}{L/min}. وعند تثبيت $Q$ عند \qty{5}{L/min} تهبط المقاومة
الكلية من 1.14 إلى $1/(1/1.14 - 1/5 + 1/0.5) =
\qty{0.37}{mmHg.s/mL}$ ويهبط الضغط إلى \qty{31}{mmHg} ---
فيُغشى على صاحبه. والجسم يرفع النتاج نحو الحالة الأولى
ويقبض الأسرّة الأخرى، فيثبت الضغط وتنال
العضلة جريانها.
\end{solution}

\begin{solution}{exo:b2:blood-pressure:11}
(أ) قرأت الكلية غير المروية ضغطًا منخفضًا، فأطلقت الرينين،
ورفع الأنجيوتنسين والألدوستيرون المقاومةَ والحجمَ.
(ب) \omterm{def:b2:blood-pressure:baroreflex}{والمنعكس الضغطي} يعزل التغيرات لكنه يعيد ضبط نفسه على أي
مستوى يدوم، ولا يستطيع تغيير احتباس الكلية للحجم.
(ج) وقد زال منبع الرينين، وزال العضو الذي يطلب ضغطًا أعلى.
(د) وكان سيحصر التحويل إلى أنجيوتنسين II: فمقاومة أقل،
وألدوستيرون أقل، وضغط أخفض.
\end{solution}

\begin{solution}{exo:b2:blood-pressure:12}
\omterm{def:b2:blood-pressure:baroreflex}{المنعكس الضغطي} يفعل في ثانية بربح نحو أربعة وبقيمة
مرجعية تنجرف إلى الضغط السائد: فهو يخمد التقلبات
ولا يستطيع إمساك مستوى. أما الكلية فتفعل على مدى ساعات وأيام،
بربح لا نهائي في الواقع --- إذ تمضي في الطرح أو الاحتباس حتى
يُبلغ الضغط الذي ضُبطت عليه --- وقيمتها المرجعية مثبتة
بفيزيولوجيتها هي: فهي تمسك المستوى. ومن ثم فالارتفاع المزمن
يعني أن قيمة الكلية المرجعية قد تحركت، ولا يخفضه إلا دواء يفعل
في حلقة الكلية أو في المقاومة التي تطلبها.
\end{solution}

\begin{solution}{pb:b2:blood-pressure:1}
\textbf{1.} $1/R = 1/7.6 + 1/22.8 + 1/4.1 + 1/5.2 + 1/5.7 + 1/19 +
1/28.5 = 0.875$: أي $R = \qty{1.14}{mmHg.s/mL}$؛ و$\bar P/Q = 95/83.3 =
1.14$.
\textbf{2.} $\bar P/R_i$: الدماغ 0.75، \omterm{def:b2:heart:anatomy}{والقلب} 0.25، والأمعاء والكبد
1.39، والكليتان 1.10، والعضلات 1.0، والجلد 0.30، وغيرها \qty{0.20}{L/min}؛
والكسور 15 و5 و28 و22 و20 و6 و\qty{4}{\%}.
\textbf{3.} $5000/70 = \qty{71}{mL}$؛ و$5\times 50 = \qty{250}{mL}$ من
الأكسجين في الدقيقة.
\textbf{4.} الدماغ $0.75/1.4 = \qty{0.54}{L/min/kg}$؛ والكليتان $1.1/0.3
= \qty{3.7}{L/min/kg}$، أي سبعة أمثال الدماغ وخمسين مثلًا لمتوسط
الجسم: فالكليتان مرويتان للترشيح لا للتغذية.
\textbf{5.} إذ يشترك كل الأعضاء في ضغط واحد، يسحب كل منها
الجريان الذي تقرره مقاومته هو؛ وتنظيم الجريانات مركزيًا كان
سيجوّع أي عضو يتغير طلبه. فالضغط هو العملة
المشتركة.
\textbf{6.} \omterm{def:b2:heart:anatomy}{القلب}: $250\times 0.2\times 0.7 = \qty{35}{mL/min}$؛
والعضلات: $1000\times 0.2\times 0.25 = \qty{50}{mL/min}$. \omterm{def:b2:heart:anatomy}{والقلب}
يستخلص أصلًا معظم ما يمر به، فلا يكسب أكسجينًا إلا
بمزيد من الجريان، وهو سبب توسع شرايينه ساعة يعمل
أشد.
\textbf{7.} هبوط $Q$ بمقدار \qty{30}{\%} عند $R$ ثابت: فيهبط $\bar P$
بمقدار \qty{30}{\%}، أي \qty{28.5}{mmHg}، إلى \qty{66.5}{mmHg}.
\textbf{8.} $28.5/5 = \qty{5.7}{mmHg}$: أي نحو \qty{89}{mmHg}.
\textbf{9.} $Q = 85\times 0.7\times 71 = \qty{4.2}{L/min} =
\qty{70}{mL/s}$؛ و$R = 89.3/70 = \qty{1.27}{mmHg.s/mL}$، أي ارتفاع
\qty{11}{\%}.
\textbf{10.} $1060\times 9.81\times 0.4 = \qty{4160}{Pa} =
\qty{31}{mmHg}$؛ والدماغ $66.5 - 31 = \qty{35}{mmHg}$ قبل المنعكس،
و$89 - 31 = \qty{58}{mmHg}$ بعده.
\textbf{11.} $28.5/2 = \qty{14}{mmHg}$: أي \qty{81}{mmHg} عند \omterm{def:b2:heart:anatomy}{القلب}،
و\qty{50}{mmHg} عند الدماغ --- فدوخة، ورؤية رمادية، وشبه غشي
عند كل نهوض من كرسي.
\textbf{12.} الاضطجاع يزيل فقد \qty{31}{mmHg} الهيدروستاتي
ويصرف الدم الراكد راجعًا إلى \omterm{def:b2:heart:anatomy}{القلب}؛ فيعود \omterm{prop:b2:heart:cycle}{حجم الضربة} وضغط
الدماغ في غضون نبضات قليلة.
\textbf{13.} $45/5 = \qty{9}{mmHg}$: أي \qty{86}{mmHg}.
\textbf{14.} $1/R = 0.875 - 1/19 - 1/4.1 + 1/38 + 1/8.2 = 0.727$: أي $R =
\qty{1.38}{mmHg.s/mL}$. والدماغ $86/7.6 = \qty{11.3}{mL/s} =
\qty{0.68}{L/min}$؛ والجلد $86/38 = \qty{2.3}{mL/s} = \qty{0.14}{L/min}$؛
والأمعاء $86/8.2 = \qty{10.5}{mL/s} = \qty{0.63}{L/min}$.
\textbf{15.} بدم أقل \omterm{def:b2:blood-circulation:vessels}{وشرينات} أضيق يهبط الضغط الشعري
تحت الضغط الجرمي على امتداد \omterm{def:b2:blood-circulation:vessels}{الشعيرة} كلها،
فيُمتص السائل بدل أن يُرشح؛ وتُمدَّد \omterm{def:b2:blood-circulation:blood}{البلازما}
فيهبط \omterm{def:b2:blood-circulation:blood}{الهيماتوكريت}.
\textbf{16.} الألدوستيرون، عبر الرينين والأنجيوتنسين، محرَّضًا بانخفاض
تروية الكلية وبالدفع الودّي؛ \omterm{prop:b2:blood-pressure:raas}{والهرمون المضاد للإبالة}،
محرَّضًا بهبوط الحجم وبارتفاع تركيز
\omterm{def:b2:blood-circulation:blood}{البلازما}.
\textbf{17.} $5\times 10^{12}$ خلية بمعدل $2.5\times 10^{6}$ في الثانية:
أي $2\times 10^{6}$ ثانية، نحو ثلاثة أسابيع.
\textbf{18.} تحت \qty{60}{mmHg} نحوًا يكون منحني المنعكس مستويًا:
فالأوعية في أضيق ما تكون \omterm{def:b2:heart:anatomy}{والقلب} في أسرع ما يكون، والربح
صفر، وكل فقد إضافي يخفض الضغط بكامل
مقداره؛ ثم تطلق الأنسجة غير المروية حمضًا وتتوسع،
فينهار الضغط.
\textbf{19.} العضلات $110/0.33 = \qty{333}{mL/s} = \qty{20}{L/min}$؛
والجلد $110/5 = \qty{22}{mL/s} = \qty{1.3}{L/min}$.
\textbf{20.} مقاومة الدماغ $8.8$ (بجريان 0.75 عند 110)، \omterm{def:b2:heart:anatomy}{والقلب} $6.6$
(بجريان 1.0): $1/R = 1/8.8 + 1/6.6 + 1/8.2 + 1/10.4 + 1/0.33 + 1/5 +
1/28.5 = 3.75$: أي $R = \qty{0.27}{mmHg.s/mL}$؛ و$Q = 110/0.27 =
\qty{410}{mL/s} = \qty{25}{L/min}$.
\textbf{21.} $\num{24700}/180 = \qty{137}{mL}$.
\textbf{22.} $24.7\times(200 - 40) = \qty{3950}{mL/min}$، أي ستة عشر مثلًا
لقيمة الراحة \qty{250}{mL/min}.
\textbf{23.} الجريان $\times 4.9$، والاستخلاص $160/50 = \times 3.2$؛
و$4.9\times 3.2 = 16$.
\textbf{24.} الأمر الآتي من القشر الحركي، والإشارات الآتية من
مستقبلات العضلات، تعيد ترجيح دخل مستقبلات الضغط في
البصلة حتى يدافع المنعكس عن \qty{110}{mmHg} بدل 95 ---
فالقيمة المرجعية مرفوعة، والحلقة غير معطّلة.
\textbf{25.} الهبوط المتبقي عند الوقوف \qty{5.7}{mmHg}؛ و\qty{86}{mmHg}
بعد النزف؛ و\qty{25}{L/min} و\qty{4}{L} من الأكسجين في
الدقيقة في الجهد.
\end{solution}
